\section{Approach}
\label{sec:approach}

Figure~\ref{fig:overview} shows the steps of \tool. The main idea is to
read a lifecycle rule from a data file, rather than hard-coding it. \tool
checks that rule against a per-function control-flow graph built from the
real parse tree of each C source file. This paper implements two such
rules against the same parser and CFG builder: Protocol B
(Section~\ref{sec:step2}), a release/reset rule derived from
CVE-2022-4450, and Protocol A (Section~\ref{sec:rq5}), a differently-shaped
establish/detach-before-release rule for the BIO-chain invariant behind
CVE-2023-0215. We return to what each rule does and does not generalize
to in Section~\ref{sec:threats}.

\begin{figure}[t]
  \centering
  %% Figure: overview of \tool (Protocol B). Steps left to right; the dashed
%% back-edge is the outer loop over pinned OpenSSL snapshots in run_scan.sh,
%% not an iteration inside one file's analysis.
\begin{tikzpicture}[
    x=1cm, y=1cm, font=\scriptsize,
    box/.style={draw, rounded corners=2pt, align=center, inner sep=2pt,
                minimum height=10mm, text width=17mm, fill=white},
    outbox/.style={draw, align=left, inner sep=2pt, fill=black!6,
                text width=22mm, font=\tiny},
    flow/.style={-{Stealth[length=1.6mm]}, line width=0.5pt},
    back/.style={-{Stealth[length=1.6mm]}, line width=0.5pt, dash pattern=on 2pt off 1.5pt},
  ]
  \node[box, text width=15mm] (in) at (-0.3,0) {C source\\+ rule (JSON)};
  \node[box] (s1) at (2.1,0) {\textbf{1} parse \&\\ build CFG};
  \node[box] (s2) at (4.3,0) {\textbf{2} find release\\ sites; search\\ paths to exit};
  \node[box] (s3) at (6.5,0) {\textbf{3} classify \&\\ extract witness};
  \node[outbox] (o) at (4.3,-1.8) {per-site result:\\ conforming /\\ violation /\\ suppressed\\ + witness path};
  \draw[flow] (in) -- (s1);
  \draw[flow] (s1) -- (s2);
  \draw[flow] (s2) -- (s3);
  \draw[flow] (s3.south) |- (o.east);
  \draw[back] (s3.north) -- ++(0,0.5) -| node[pos=0.3, above] {next snapshot} (s1.north);
\end{tikzpicture}

  \caption{\tool takes a C source tree and a rule stored as data
    (Section~\ref{sec:rule}). It parses every file once, finds every site
    that releases a tracked pointer, and searches forward in the
    control-flow graph for a path to function exit that never resets it.
    Each flagged site keeps a witness path. The outer loop (dashed) reruns
    this per pinned OpenSSL snapshot; nothing loops within one file's
    analysis.}
  \label{fig:overview}
\end{figure}

\subsection{The rule, as data}
\label{sec:rule}
The rule for Protocol B is a JSON file, not a code path: it names no
OpenSSL function. It tracks parameters of pointer-to-pointer type; a
\emph{release} is any call whose name matches the regular expression
\texttt{(?i)free} and whose argument is \texttt{*P}; a \emph{reset} is an
assignment to \texttt{*P}, or, by a configurable flag, passing \texttt{P}
itself to a later call. The invariant is: if a function releases
\texttt{*P}, every path from that release to function exit must reach a
reset first. We do not claim the idea of an output-parameter reset
invariant as new; what we built is a parser and path-search engine that
checks it against real source, rather than against one patch's text.

\subsection{Step 1: Parsing and the control-flow graph}
\label{sec:step1}
\tool parses each C file with \texttt{tree-sitter} and
\texttt{tree-sitter-c}. We do not claim parsing itself as new; everything
from here on is. For each function, we build a control-flow graph whose
nodes are simple statements, branch and loop conditions, labels, and join
points; \texttt{goto} targets, including error labels such as
\texttt{err:}, become real edges, so an error path that jumps past a reset
is visible to the path search in Step~2. Three approximations keep the
graph tractable and are stated as threats to validity
(Section~\ref{sec:threats}): a compound condition joined by \texttt{\&\&}
or \texttt{||} is one node, without separate short-circuit edges; a
\texttt{\#if}/\texttt{\#ifdef} block is treated as a branch, keeping both
variants rather than resolving the macro; and macro bodies are not
expanded, so a macro call is treated as an ordinary call expression.

\subsection{Step 2: Release sites and the path search}
\label{sec:step2}
For every parameter \texttt{P} of pointer-to-pointer type or deeper, \tool
scans each control-flow node for calls matching the release pattern and
for assignments matching the reset pattern. From every release site, it
searches the control-flow graph forward, breadth-first, for a path to
function exit. The search stops following a branch once it crosses a
reset, and the release's own dereference of \texttt{*P} licenses one
pruning rule: because the release itself has already dereferenced
\texttt{*P}, \texttt{P} is non-null on every path after it, so a later
branch that is live only when \texttt{P == NULL} is dead code and is not
explored. If the search reaches function exit without crossing a reset,
the site is a candidate violation and the path taken becomes its witness.

Two refinements narrow this candidate set. We disclose both as chosen
after reading real output, not derived from the rule alone. \emph{Deref-
after-release suppression} drops a site where \texttt{*P} is read again
before any reset. This signals a partial release and reuse, not a
dangling pointer: for example, a function that frees only part of a
struct's contents and then reinitializes it. The second refinement is the
\texttt{NULL}-guard pruning described above, added after it produced one
infeasible-path false positive on a whole-tree scan. Both are unit-tested
(Section~\ref{sec:eval}) and are listed, not hidden, in each result.

\subsection{Step 3: Classification and witness extraction}
\label{sec:step3}
Each release site is labeled \emph{conforming} if every path from it to
function exit passes a reset; \emph{violation} if some path reaches exit
without one; or \emph{suppressed} if the path rule would flag it but one
of the two refinements in Step~2 applies. A suppressed site is reported,
with its reason, alongside conforming and violation sites; it is never
dropped silently. For a violation, \tool records the source lines on the
witness path. A reader can then check the finding against the source
without re-running the tool. \texttt{compare.py} diffs two runs of the
same files at different revisions and reports exactly which sites a given
commit removed or added. Section~\ref{sec:eval} uses this to measure
whether a real fix changes only the sites its own diff touches.
